\nadpis {The Cookbook of Matfyz}
\noindent Whether you are reading this on your laptop, your phone or you are holding the printed version of this not just ordinary Cookbook, welcome. Hopefully you will find it to be quite useful.


Unfortunately this Cookbook will not provide you with a lamb stew recipe because it is too complicated.
However, you will find here many recipes describing how to survive studying at Matfyz (which in Czech is short for ``the Faculty of Mathematics and Physics'').
Although we would like to let you discover the beauties of our faculty buildings, metro tunnels and Prague's numerous spires on your own, we believe that a little help might come in handy, at least in the beginning.
It might spare you a lot of embarrassing situations and misunderstandings and it will certainly make your transition to a new school less stressful.
This Cookbook was created as a student handbook giving details about the buildings of MFF (or MFF UK which in Czech stands for ``the Faculty of Mathematics and Physics at Charles University''), life in dormitories (mostly in the student residence 17. listopadu) and using the Prague public transport.


The first Student Cookbook was created in 1994 thanks to Martin Krynický (the 1998 version is available here).
From that time on new and new authors, members of Spolek Matfyzák (a group of the MFF's students organising various events and helping other students), try to improve the Cookbook by correcting old mistakes and subsequently adding new ones.
It is freely available online (\url{kucharka.matfyzak.cz}(CZ)) since 2011.


We will describe and explain almost everything here but some things are impossible to capture: the atmosphere of lectures, the presentations followed by applause or accompanied with laugh or loud snoring, the sleepless nights, the weeks and semesters spent writing the physics laboratory protocols, the first grades in SIS (Study Information System), the exams successfully passed talking over tea offered by a professor or even the failed exams.